\section{Diskussion}
\label{sec:Diskussion}

Das durchgeführte Experiment zum BB84-Protokoll mittels Thorlabs Komponenten kann 
als sehr erfolgreich gewertet werden. Für die Schlüsselgenerierung aus \autoref{sec:a1}
ergibt sich eine Fehlerrate von 0\%. Dieser experimentelle Befund zeugt von einem 
sehr präzisen Aufbau und geeigneten Umgebungsfaktoren. Gleiches ergibt sich in 
\autoref{sec:a2}. In diesem Teil geht es um die Hinzunahme des Abhörmoduls Eve. 
Eve und Bob könnten diese Fehlerquote als Indiz für Eves Abhörversuch nutzen. Dafür
würden sie in der Praxis einen Teil ihres Schlüssels vergleichen, um Fehler 
aufzudecken. Falls Fehler zu verzeichnen sind, ist Eve aufgedeckt und der Schlüssel 
sollte von Alice und Bob nicht verwendet werden. Hier wurden alle Bits des Schlüssels
verglichen, da es sonst zu wenig Vergleichsmaterial gegeben hätte. Der ermittelte 
Wert entspricht somit der in \autoref{sec:bb84} erläuterten Theorie.
\vspace{0.5em}
\\
\noindent Im zweiten Teil des Experiments (siehe \autoref{sec:a3} und \autoref{sec:a3}) werden Messungen auf Basis eines ternären
Systems mit 2 Quellen unterschiedlicher Wellenlängen wiederholt. Das geschieht mit 
und ohne Abhörmodul. Wie zuvor ergibt sich die Fehlerrate ohne Eve zu 0\%. Bei 
Hinzunahme belaufen sich die Error Rate auf ganze 42\%, die Decoy Rate auf 21\%.
Die Error Rate überschreitet den Erwartungswert von 25\% deutlich. Diese Abweichung 
wird durch mehrere Faktoren zu erklären sein. Im zweiten Teil des Experiments kam 
es bei der Justage zu Problemen, wo die Messdetektoren auch nach mehrfacher Korrektur 
nicht zu den Ergebnissen führten, wie vom Kontrollblatt \autoref{fig:d3} angegeben.
Es könnte sich ein systematischer Fehler gebildet haben, der nach finaler Justage 
während des Messvorgangs erneut aufgetaucht ist und die Messrate dementsprechend in 
die Höhe getrieben hat. Darüber hinaus ist die Statistik mit 33 ausgewerteten Bits 
relativ klein, sodass zufällige Schwankungen einen großen Einfluss haben können. 
Eine größere Anzahl an gesendeten Pulsen würde das Ergebnis näher an den theoretischen
Erwartungswert bringen. Besonders aufschlussreich ist die Decoy Rate von 21\%. Da
von Alice aus diese Zustände bewusst nicht gesendet werden, ist ihr Auftreten ein 
direkter Nachweis für eine Störung des Quantenkanals. Im hier durchgeführten 
Experiment ist diese Störung eindeutig auf Eves Eingriff zurückzuführen. Die
Detektion von Decoy States bietet somit einen klaren und einfach zu interpretierenden
Indikator für Abhörversuche. Ein Vorteil gegenüber dem einfachen BB84-Protokoll,
das nur auf die Fehlerrate angewiesen ist.
\\
\noindent Zusammenfassend lässt sich sagen, dass der Versuch die Funktionsweise des BB84-Protokolls
sowie die Detektion von Abhörversuchen mittels Fehlerrate und Decoy States erfolgreich
demonstriert.